\documentclass[11pt,a4paper]{article}
\usepackage[utf8]{inputenc}
\usepackage[T1]{fontenc}
\usepackage[spanish]{babel}
\usepackage{geometry}
\geometry{margin=2.5cm}
\usepackage{hyperref}
\usepackage{longtable}
\usepackage{array}
\usepackage{booktabs}
\usepackage{enumitem}
\usepackage{verbatim}
\usepackage{listings}
\usepackage{graphicx}

\title{SEGUNDA EVALUACIÓN DE PROYECTO\\Análisis Geoespacial Aplicado\\Presentación Final e Informe Técnico}
\author{Prof.\ Francisco Parra O.\\\href{mailto:francisco.parra@usach.cl}{francisco.parra@usach.cl}}
\date{Fecha de publicación: 9 de enero de 2026\\Fecha de presentaciones: 22 de enero de 2026\\Fecha de entrega del informe: Primera semana de marzo de 2026}

\begin{document}
\maketitle
\tableofcontents
\newpage

\section{Descripción General}
Esta segunda evaluación constituye el hito final del proyecto semestral. A diferencia de la primera
entrega (que evaluaba factibilidad y avances), esta evaluación se centra en:
\begin{itemize}[leftmargin=*]
\item \textbf{Resultados completos:} Análisis terminados con conclusiones sólidas
\item \textbf{Implementación funcional:} Código completo, reproducible y documentado
\item \textbf{Producto final:} Dashboard, aplicación o sistema demostrable
\item \textbf{Contribución:} Impacto y valor del trabajo realizado
\item \textbf{Comunicación:} Capacidad de presentar y defender el trabajo
\end{itemize}

\subsection*{Importante}
Esta es la evaluación final del proyecto. Se espera un trabajo terminado, no un avance. Los
criterios son más exigentes que en la primera entrega.

\subsection{Diferencias con la Primera Entrega}
\begin{center}
\begin{tabular}{p{0.47\textwidth} p{0.47\textwidth}}
\toprule
\textbf{Primera Entrega (PEP 1)} & \textbf{Segunda Entrega (PEP 2)}\\
\midrule
Análisis exploratorio & Análisis completo y profundo\\
Resultados preliminares & Resultados finales con validación\\
Plan de trabajo & Trabajo ejecutado y documentado\\
Cronograma futuro & Reflexión sobre lo realizado\\
40--50\% de avance & 100\% completado\\
Verificar factibilidad & Demostrar impacto y contribución\\
Propuesta metodológica & Metodología implementada\\
Prototipo de visualización & Dashboard/aplicación funcional\\
\bottomrule
\end{tabular}
\end{center}

\section{Objetivos de la Evaluación}
\subsection{Objetivos Principales}
\begin{enumerate}[leftmargin=*]
\item Evaluar completitud: Verificar que el proyecto está terminado según lo propuesto
\item Validar resultados: Confirmar que los análisis son correctos y las conclusiones fundamentadas
\item Medir calidad técnica: Evaluar la implementación, código y reproducibilidad
\item Evaluar comunicación: Capacidad de presentar y defender el trabajo
\item Identificar contribución: Valor agregado del proyecto al problema planteado
\end{enumerate}

\section{Componentes de la Evaluación}
\subsection{Distribución de Ponderaciones}
\begin{center}
\begin{tabular}{p{0.46\textwidth} p{0.12\textwidth} p{0.36\textwidth}}
\toprule
\textbf{Componente} & \textbf{Peso} & \textbf{Descripción}\\
\midrule
Informe Técnico Final & 35\% & Documento completo con resultados y conclusiones\\
Presentación Oral & 30\% & Exposición de 20 min + 10 de preguntas\\
Código y Reproducibilidad & 20\% & Repositorio completo y funcional\\
Producto Final (Dashboard/App) & 10\% & Aplicación interactiva demostrable\\
Trabajo en Equipo & 5\% & Evidencia de colaboración efectiva\\
\midrule
Total & 100\% & \\
\bottomrule
\end{tabular}
\end{center}

\subsection*{Recomendación}
La presentación tiene mayor peso (30\%) porque es la instancia donde demuestran dominio
del tema y capacidad de comunicar resultados. Prepárense bien.

\section{Requisitos por Componente}
\subsection{A.\ Informe Técnico Final (35\%)}
\subsubsection*{Requisitos del Informe}
\begin{itemize}[leftmargin=*]
\item Extensión: 25--35 páginas (sin anexos)
\item Formato: \LaTeX\ o Markdown convertido a PDF
\item Entrega: Primera semana de marzo de 2026
\end{itemize}

\subsubsection{Estructura Obligatoria del Informe}
\begin{enumerate}[leftmargin=*]
\item \textbf{Resumen Ejecutivo (1 página)}
  \begin{itemize}[leftmargin=*]
  \item Problema abordado y motivación
  \item Metodología aplicada
  \item Principales resultados y hallazgos
  \item Conclusiones clave
  \item Impacto y contribución
  \end{itemize}
\item \textbf{Introducción (2--3 páginas)}
  \begin{itemize}[leftmargin=*]
  \item Contexto actualizado del problema
  \item Objetivos (cumplidos vs.\ propuestos)
  \item Preguntas de investigación respondidas
  \item Alcance final del proyecto
  \end{itemize}
\item \textbf{Marco Teórico (3--4 páginas)}
  \begin{itemize}[leftmargin=*]
  \item Revisión de literatura (mínimo 15 referencias)
  \item Conceptos geoespaciales aplicados
  \item Justificación teórica de métodos utilizados
  \end{itemize}
\item \textbf{Área de Estudio (2 páginas)}
  \begin{itemize}[leftmargin=*]
  \item Delimitación geográfica final
  \item Características relevantes para el análisis
  \item Mapas de contexto de alta calidad
  \end{itemize}
\item \textbf{Datos y Metodología (5--6 páginas)}
  \begin{itemize}[leftmargin=*]
  \item Fuentes de datos utilizadas (tabla detallada)
  \item Procesamiento aplicado (con justificación)
  \item Diagrama de flujo metodológico completo
  \item Métodos de análisis espacial empleados
  \item Herramientas y librerías utilizadas
  \item Limitaciones metodológicas
  \end{itemize}
\item \textbf{Resultados (6--8 páginas) --- SECCIÓN CLAVE}
  \begin{itemize}[leftmargin=*]
  \item Presentación sistemática de hallazgos
  \item Visualizaciones de alta calidad
  \item Análisis estadístico espacial completo
  \item Respuesta a cada pregunta de investigación
  \item Validación de resultados (si aplica)
  \end{itemize}
\item \textbf{Discusión (3--4 páginas) --- NUEVA SECCIÓN}
  \begin{itemize}[leftmargin=*]
  \item Interpretación de resultados
  \item Comparación con literatura/casos similares
  \item Implicancias prácticas de los hallazgos
  \item Limitaciones del estudio
  \item Trabajo futuro propuesto
  \end{itemize}
\item \textbf{Conclusiones (2 páginas)}
  \begin{itemize}[leftmargin=*]
  \item Síntesis de principales hallazgos
  \item Cumplimiento de objetivos
  \item Contribución del proyecto
  \item Recomendaciones concretas
  \end{itemize}
\item \textbf{Referencias} (mínimo 15)
\item \textbf{Anexos}
  \begin{itemize}[leftmargin=*]
  \item Código relevante comentado
  \item Mapas adicionales
  \item Tablas de datos extensas
  \item Manual de usuario del dashboard (si aplica)
  \end{itemize}
\end{enumerate}

\subsection{B.\ Presentación Oral (30\%)}
\subsubsection*{Formato de Presentación}
\begin{itemize}[leftmargin=*]
\item Duración: 20 minutos de exposición + 10 minutos de preguntas
\item Formato: Presencial con apoyo de slides
\item Fecha: 22 de enero de 2026
\item Participación: Todos los integrantes deben exponer
\end{itemize}

\subsubsection{Estructura Sugerida de la Presentación}
\begin{enumerate}[leftmargin=*]
\item \textbf{Introducción y Motivación (3 min)}
  \begin{itemize}[leftmargin=*]
  \item Problema y por qué es relevante
  \item Objetivos del proyecto
  \end{itemize}
\item \textbf{Datos y Metodología (4 min)}
  \begin{itemize}[leftmargin=*]
  \item Fuentes de datos (mostrar ejemplos visuales)
  \item Pipeline de procesamiento
  \item Métodos de análisis aplicados
  \end{itemize}
\item \textbf{Resultados Principales (8 min) --- LA PARTE MÁS IMPORTANTE}
  \begin{itemize}[leftmargin=*]
  \item Visualizaciones clave
  \item Hallazgos principales
  \item Demo del dashboard/aplicación
  \end{itemize}
\item \textbf{Conclusiones y Discusión (3 min)}
  \begin{itemize}[leftmargin=*]
  \item Qué aprendimos
  \item Implicancias prácticas
  \item Limitaciones y trabajo futuro
  \end{itemize}
\item \textbf{Cierre (2 min)}
  \begin{itemize}[leftmargin=*]
  \item Contribución del proyecto
  \item Reflexión del equipo
  \end{itemize}
\end{enumerate}

\subsubsection*{Importante}
La demo del dashboard/aplicación es obligatoria. Debe funcionar en vivo. Tengan un plan
B (video grabado) por si hay problemas técnicos.

\subsection{C.\ Código y Reproducibilidad (20\%)}
\subsubsection{Estructura del Repositorio}
\begin{lstlisting}[basicstyle=\ttfamily\small]
proyecto/
README.md # Documentacion completa
requirements.txt # Dependencias Python
environment.yml # Alternativa conda (opcional)
data/
raw/ # Datos originales o scripts de descarga
processed/ # Datos procesados
notebooks/
01_data_acquisition.ipynb
02_preprocessing.ipynb
03_exploratory_analysis.ipynb
04_spatial_analysis.ipynb
05_visualization.ipynb
src/
__init__.py
data_loader.py # Funciones de carga de datos
preprocessing.py # Funciones de procesamiento
analysis.py # Funciones de analisis
visualization.py # Funciones de visualizacion
utils.py # Utilidades generales
app/
streamlit_app.py # Dashboard principal
pages/ # Paginas adicionales (si aplica)
outputs/
figures/ # Graficos generados
maps/ # Mapas producidos
reports/ # Reportes generados
docs/
informe_final.pdf
presentacion.pdf
\end{lstlisting}

\subsubsection{Criterios de Calidad del Código}
\begin{itemize}[leftmargin=*]
\item Documentación: Cada función con docstring, README completo
\item Modularidad: Código organizado en módulos reutilizables
\item Reproducibilidad: Instrucciones claras paso a paso
\item Manejo de datos: Scripts de descarga si datos $>$ 100MB
\item Versionamiento: Commits descriptivos de todos los integrantes
\item Estilo: Código limpio siguiendo PEP 8
\end{itemize}

\subsection{D.\ Producto Final: Dashboard/Aplicación (10\%)}
\subsubsection*{Requisitos del Dashboard}
Se espera una aplicación interactiva que permita explorar los resultados del proyecto. Puede
ser:
\begin{itemize}[leftmargin=*]
\item Streamlit (recomendado)
\item Panel/Voila
\item Dash (Plotly)
\item Aplicación web estática con Folium
\end{itemize}

\subsubsection{Funcionalidades Mínimas}
\begin{enumerate}[leftmargin=*]
\item Visualización de mapa interactivo con los datos del proyecto
\item Filtros dinámicos para explorar subconjuntos de datos
\item Al menos 2 gráficos que se actualicen con los filtros
\item Información contextual sobre el proyecto
\item Instrucción de uso básica
\end{enumerate}

\subsubsection*{Recomendación}
Si desplegaron el dashboard en Streamlit Cloud u otro servicio, incluyan el enlace en el
README. Esto demuestra profesionalismo y facilita la evaluación.

\subsection{E.\ Trabajo en Equipo (5\%)}
Se evaluará:
\begin{itemize}[leftmargin=*]
\item Commits de todos los integrantes
\item Distribución equitativa del trabajo
\item Evidencia de colaboración (issues, PRs, comentarios)
\item Participación equitativa en la presentación
\end{itemize}

\section{Rúbrica de Evaluación Detallada}
\subsection{Niveles de Desempeño}
\begin{longtable}{p{0.18\textwidth} p{0.20\textwidth} p{0.20\textwidth} p{0.20\textwidth} p{0.18\textwidth}}
\toprule
\textbf{Criterio} & \textbf{Excelente (7.0)} & \textbf{Bueno (6.0)} & \textbf{Suficiente (5.0)} & \textbf{Insuficiente ($<4.0$)}\\
\midrule
\endfirsthead
\toprule
\textbf{Criterio} & \textbf{Excelente (7.0)} & \textbf{Bueno (6.0)} & \textbf{Suficiente (5.0)} & \textbf{Insuficiente ($<4.0$)}\\
\midrule
\endhead
Resultados y Análisis &
Análisis sofisticado con múltiples métodos espaciales. Hallazgos profundos y bien fundamentados. Validación rigurosa. &
Buen análisis con métodos apropiados. Hallazgos claros y fundamentados. &
Análisis básico pero correcto. Hallazgos limitados pero válidos. &
Análisis superficial, incorrecto o ausente. Sin hallazgos claros.\\
\midrule
Visualizaciones &
Excepcionales: claras, estéticas, interactivas. Comunican efectivamente. Mapas de calidad profesional. &
Buena calidad, claras y bien diseñadas. Mapas correctos con elementos cartográficos. &
Correctas pero básicas. Mapas funcionales sin refinamiento. &
Pobres, confusas o con errores. Mapas de baja calidad.\\
\midrule
Dashboard/App &
Funcional, intuitivo, estético. Múltiples funcionalidades. Desplegado en la nube. &
Funcional y útil. Cumple requisitos mínimos. Bien diseñado. &
Básico pero funciona. Interactividad limitada. &
No funciona, muy básico o ausente.\\
\midrule
Código &
Limpio, modular, bien documentado. Completamente reproducible. Tests incluidos. &
Bien estructurado y funcional. Documentación adecuada. &
Funciona pero poco organizado. Documentación mínima. &
Desorganizado, sin documentar o no funciona.\\
\midrule
Informe &
Completo, bien redactado, estructura impecable. Discusión profunda. Referencias de calidad. &
Bien estructurado, redacción clara. Contenido completo. &
Estructura correcta, contenido básico. Algunas secciones débiles. &
Incompleto, mal estructurado o mal redactado.\\
\midrule
Presentación &
Clara, fluida, dominio total del tema. Responde preguntas con profundidad. Demo impecable. &
Buena presentación, responde bien. Demo funciona correctamente. &
Presentación correcta. Demo con problemas menores. &
Confusa, incompleta. Demo falla o ausente.\\
\midrule
Conclusiones &
Profundas, bien fundamentadas. Identifican contribución clara. Proponen trabajo futuro concreto. &
Claras y pertinentes. Contribución identificada. &
Básicas pero correctas. Contribución vaga. &
Superficiales, incorrectas o ausentes.\\
\bottomrule
\end{longtable}

\section{Checklist de Entregables}
\subsection*{Importante}
Todos los items son obligatorios. La ausencia de cualquiera resultará en descuento.
\begin{center}
\begin{tabular}{p{0.70\textwidth} p{0.20\textwidth}}
\toprule
\textbf{Item} & \textbf{Estado}\\
\midrule
Informe final en PDF (25--35 páginas) & Pendiente\\
Presentación en PDF o PPTX & Pendiente\\
Repositorio GitHub completo y documentado & Pendiente\\
README con instrucciones de reproducción & Pendiente\\
requirements.txt o environment.yml & Pendiente\\
Dashboard/aplicación funcional & Pendiente\\
Mínimo 5 mapas temáticos de alta calidad & Pendiente\\
Mínimo 8 gráficos estadísticos & Pendiente\\
Notebook de análisis completo y ejecutable & Pendiente\\
Código modular en carpeta src/ & Pendiente\\
Marco teórico con mínimo 15 referencias & Pendiente\\
Sección de discusión completa & Pendiente\\
Conclusiones con contribución identificada & Pendiente\\
Commits de todos los integrantes & Pendiente\\
Demo funcional para presentación & Pendiente\\
\bottomrule
\end{tabular}
\end{center}

\section{Criterios de Penalización}
\subsection*{Importante}
\begin{itemize}[leftmargin=*]
\item Entrega tardía del informe: -1.0 punto por día de atraso
\item Falta de integrante a presentación: -0.5 puntos (salvo justificación médica)
\item Código no reproducible: -0.5 puntos
\item Dashboard no funciona: -0.5 puntos
\item Sin commits de algún integrante: -0.3 puntos
\item Plagio o copia: Nota 1.0 y revisión por comité de ética
\end{itemize}

\section{Sesiones Prácticas de Cierre}
Se realizarán 3 sesiones prácticas online antes de la presentación final:
\begin{center}
\begin{tabular}{p{0.35\textwidth} p{0.18\textwidth} p{0.37\textwidth}}
\toprule
\textbf{Fecha} & \textbf{Duración} & \textbf{Foco}\\
\midrule
Lunes 13 de enero & 1:20 hrs & Completar análisis y resultados\\
Miércoles 15 de enero & 1:20 hrs & Visualizaciones y dashboard\\
Lunes 20 de enero & 1:20 hrs & Preparación de presentación\\
Jueves 22 de enero & --- & Presentaciones (presencial)\\
\bottomrule
\end{tabular}
\end{center}

\subsection*{Importante}
Es obligación de cada grupo tener al menos 2 integrantes presentes en cada sesión práctica.
La falta no justificada resultará en -0.3 puntos en la nota final.

\section{Fechas Importantes}
\subsection*{Calendario}
\begin{itemize}[leftmargin=*]
\item Publicación de esta guía: 9 de enero de 2026
\item Sesión práctica 1: 13 de enero de 2026
\item Sesión práctica 2: 15 de enero de 2026
\item Sesión práctica 3: 20 de enero de 2026
\item Presentaciones finales: 22 de enero de 2026 (presencial)
\item Entrega informe final: Primera semana de marzo de 2026
\end{itemize}

\section{Contacto}
\begin{itemize}[leftmargin=*]
\item Email: \href{mailto:francisco.parra@usach.cl}{francisco.parra@usach.cl}
\item Consultas: Durante las sesiones prácticas o por email
\end{itemize}

\bigskip
\noindent Éxito en su evaluación final!
\end{document}
